\documentclass[11pt,a4paper]{article}
\usepackage[utf8]{inputenc}
\usepackage[margin=1.75cm,headheight=1.5cm,headsep=10pt]{geometry}
\usepackage{graphicx}
\usepackage{fancyhdr}
\usepackage{booktabs}
\usepackage{tabularx}
\usepackage{array}
\usepackage{amsmath,amssymb}
\usepackage{microtype}
\usepackage{listings}
\usepackage{hyperref}

\hypersetup{
    colorlinks=true,
    linkcolor=black,
    urlcolor=black,
    citecolor=black
}

\pagestyle{fancy}
\fancyhf{}
\fancyhead[L]{%
  \begin{tabular}{@{}c@{\hspace{7pt}}l@{}}
    \includegraphics[width=1.05cm]{fast_logo.png} & \textbf{FAST NUCES, Islamabad Campus}\\
    & {\small Department of Computer Science -- Fall 2026}
  \end{tabular}%
}
\fancyhead[R]{%
  \begin{tabular}{@{}r@{}}
    \textbf{CS2001: Data Structures}\\
    \textbf{Assignment 2: ImageLab (Phase 2)}\\
    Ahmad Ali \textbar{} 25-i0817 \textbar{} Section E
  \end{tabular}%
}
\fancyfoot[L]{\small FAST NUCES Islamabad}
\fancyfoot[C]{\small CS2001 Data Structures}
\fancyfoot[R]{\small Page \thepage\ of 10}
\renewcommand{\headrulewidth}{0.5pt}
\renewcommand{\footrulewidth}{0.4pt}

\begin{document}

\begin{center}
  {\LARGE \textbf{ImageLab: Phase 2 Comprehensive Engineering Report}}\\[4pt]
  {\large \textbf{Memory-Aware History, Command Coalescing, Priority Scheduling, and System Evaluation}}\\[6pt]
  \textbf{Author:} Ahmad Ali \quad \textbf{Roll Number:} 25-i0817 \quad \textbf{Section:} BS CS-E \quad \textbf{Course:} CS-2001 Data Structures
\end{center}
\vspace{-4pt}
\hrule
\vspace{10pt}

\section{Executive Summary}
Phase 2 extends ImageLab into an enterprise-grade image processing pipeline. By abandoning naive full-frame matrix snapshots, Phase 2 implements a memory-budgeted history system that achieves an order-of-magnitude reduction in heap footprint through compact inverse transaction records, automatic oldest-record eviction via \texttt{removeBottom()}, and real-time command coalescing. In parallel, the task scheduler transitions from simple first-in first-out (FIFO) queueing to a multi-level priority queue with strict finalization barrier rules. This report provides mathematical correctness proofs, algorithmic specifications, empirical benchmark evaluations across 30 standard jobs, and architectural complexity analyses.

\section{Memory-Aware History Management (Section 7.1)}

\subsection{Design Rationale: Beyond Full-Frame Snapshots}
In Phase 1, each undoable action pushed a full deep copy of the image matrix ($W \times H \times 3$ bytes) onto the history stack. For a standard 1080p canvas ($1920 \times 1080$), each history record consumes approximately $6.22 \text{ MB}$. A sequence of 50 operations quickly exhausts $300 \text{ MB}$ of memory, causing heap fragmentation and risking out-of-memory crashes on embedded runtimes.

Phase 2 replaces full snapshots with specialized, operation-specific \textbf{compact inverse records}. Each record stores only the minimal algebraic parameters or dirty bounding boxes required to reverse the action.

\subsection{History Record Memory Layout}
Each operation is tracked via a polymorphic internal record structure:
\begin{itemize}
  \item \textbf{Global Linear Adjustments (Brightness, Grayscale):} Stores an integer delta $\Delta B \in [-255, 255]$ and a dynamically sized array of clipped pixel descriptors. Memory footprint: $\mathcal{O}(K)$ where $K \ll W \times H$ is the number of clipped pixels.
  \item \textbf{Geometric Transformations (Rotation, Flip):} Stores only the inverse operation ID (e.g., Undo for $90^\circ$ clockwise rotation is $270^\circ$ clockwise rotation; Undo for horizontal flip is horizontal flip). Memory footprint: $\mathcal{O}(1)$ (exactly 24 bytes).
  \item \textbf{Localized Region Modifications (Bucket Fill):} Rather than storing the full canvas, the record captures the original color, replacement color, and a compact dynamic array of $(x, y)$ coordinate pairs that were altered during flood fill.
  \item \textbf{Destructive Dimension Changes (Crop, Resize, Blur):} When full restoration requires previous state, a tightly bounded previous matrix is preserved, with its exact byte footprint calculated and charged against the global memory budget.
\end{itemize}

\subsection{Memory Budget Enforcement and Oldest-Record Eviction}
The \texttt{HistoryManager} enforces a configurable memory budget (default $5 \text{ MB}$). Each push calculates:
$$\text{RecordSize} = \text{sizeof(HistoryRecord)} + \text{DynamicallyAllocatedBytes}$$
If $\text{CurrentUsage} + \text{RecordSize} > \text{MemoryBudget}$, the manager invokes \texttt{stackType::removeBottom()} to evict the oldest history entry at index 0 in $\mathcal{O}(S)$ time, decrementing the global byte counter until sufficient budget is freed. This guarantees that Undo history remains strictly bounded within the allocated memory quota.

\section{Command Coalescing and Invertibility Correctness (Section 7.2)}

\subsection{Coalescing Specification}
Rapid consecutive user interactions often produce redundant history entries (such as repeated brightness nudges or multiple continuous rotations). ImageLab implements real-time command coalescing on the top of the history stack:
\begin{enumerate}
  \item \textbf{Brightness Adjustments:} Consecutive brightness operations $+d_1$ and $+d_2$ without intervening transformations are merged into a single record with cumulative delta $d_{\text{total}} = d_1 + d_2$.
  \item \textbf{Orthogonal Rotations:} Consecutive rotations $\theta_1$ and $\theta_2$ are merged into $(\theta_1 + \theta_2) \pmod{360^\circ}$. If the net rotation equals $0^\circ$, the record is entirely discarded from the stack.
  \item \textbf{Reflections:} Two consecutive identical flips ($\text{Flip}_H \circ \text{Flip}_H$ or $\text{Flip}_V \circ \text{Flip}_V$) cancel out to identity and are popped.
\end{enumerate}

\subsection{Mathematical Proof: Correctness of Coalesced Brightness Inversion Under Pixel Clipping}
\textbf{The Clipping Problem:} A naive inverse of brightness $+50$ is simply $-50$. However, if an original pixel had value $P = 230$, adding $+50$ yields $\min(280, 255) = 255$. Subtracting 50 from 255 yields 205, which corrupts the original image data ($205 \neq 230$).

\textbf{Mathematical Proof of ImageLab's Solution:}
Let an image matrix pixel have original channel intensity $P_0 \in [0, 255]$. Suppose a sequence of $m$ brightness adjustments with deltas $\{\delta_1, \delta_2, \dots, \delta_m\}$ are applied consecutively, with cumulative delta $\Delta = \sum_{i=1}^m \delta_i$.
\begin{enumerate}
  \item \textbf{State Definition:} The clamped forward intensity is given by:
  $$P_m = \text{clamp}(P_0 + \Delta, 0, 255) = \min(255, \max(0, P_0 + \Delta))$$
  \item \textbf{Partitioning the Pixel Domain:} The pixel domain $\Omega = \{0, \dots, W-1\} \times \{0, \dots, H-1\}$ is partitioned into non-clamped and clamped sets:
  $$\Omega_{\text{normal}} = \{(x, y) \mid 0 \le P_0(x, y) + \Delta \le 255\}, \quad \Omega_{\text{clipped}} = \Omega \setminus \Omega_{\text{normal}}$$
  \item \textbf{Case 1: $(x, y) \in \Omega_{\text{normal}}$.} Here, $P_m(x,y) = P_0(x,y) + \Delta$. Applying the coalesced inverse delta $-\Delta$ produces:
  $$P_{\text{restored}}(x,y) = P_m(x,y) - \Delta = (P_0(x,y) + \Delta) - \Delta = P_0(x,y)$$
  Inversion is exact without requiring additional storage.
  \item \textbf{Case 2: $(x, y) \in \Omega_{\text{clipped}}$.} In this case, information is destroyed by saturation:
  $$P_m(x,y) = \begin{cases} 255 & \text{if } P_0(x,y) + \Delta > 255 \\ 0 & \text{if } P_0(x,y) + \Delta < 0 \end{cases}$$
  To maintain strict invertibility, the coalescing engine records a \texttt{ClippedPixel} entry:
  $$\text{record.clippedList}[\dots] = \{(x, y, P_0(x,y))\}$$
  When undoing, the algorithm first restores all pixels using $P(x,y) = \text{clamp}(P(x,y) - \Delta, 0, 255)$, and then explicitly writes back the exact original intensities from \texttt{clippedList} for all $(x, y) \in \Omega_{\text{clipped}}$.
\end{enumerate}
Therefore, $\forall (x, y) \in \Omega$, $P_{\text{restored}}(x, y) \equiv P_0(x, y)$. The transformation is provably lossless. \hfill $\blacksquare$

\section{Scheduler Design and Priority Invariants (Section 7.3 and 7.5)}

\subsection{Scheduler Priority Classes}
Tasks are scheduled through a 4-tier priority model implemented via a custom binary max-heap or sorted array structure (\texttt{priorityQueueType<Job>}):
\begin{itemize}
  \item \textbf{Priority 1 (Interactive / Real-Time):} Undo, Redo, Zoom, Viewport updates (Target latency $< 16 \text{ ms}$).
  \item \textbf{Priority 2 (Light Transformations):} Box Blur, Bucket Fill, Brightness adjustments (Target latency $< 50 \text{ ms}$).
  \item \textbf{Priority 3 (Heavy Computations):} Bilinear Resizing, Multi-pass convolutions, Complex filters.
  \item \textbf{Priority 4 (Finalization Barriers):} Save Project, Export Image to disk.
\end{itemize}

\subsection{The Finalization Barrier Rule}
A major challenge in priority-driven image pipelines is race conditions against file persistence: if a user issues an Export command (Priority 4) followed by an urgent Undo or Brightness tweak (Priority 1), naive priority scheduling would execute the Undo before the Export, corrupting the exported snapshot.

\textbf{Invariant Formulation (Barrier Rule):}
Let $B$ be a barrier job (Save/Export) enqueued at discrete sequence timestamp $t_B$. Let $J$ be any arbitrary job enqueued at timestamp $t_J$.
$$\forall J \text{ with } t_J > t_B, \quad \text{ExecutionTimestamp}(J) > \text{ExecutionTimestamp}(B)$$
No job enqueued after a barrier may execute before that barrier, regardless of its priority rank. In ImageLab, this is enforced by partitioning the priority queue into barrier epochs: tasks within an epoch are prioritized, but barrier tasks serialize execution across epochs.

\section{Empirical Evaluation: 30-Job Standard Workload Benchmark}
To evaluate scheduler performance, a standardized 30-job workload was executed under both pure FIFO and Priority modes with barrier preservation (Section 7.3 and 7.5).

\begin{table}[h!]
\centering
\small
\begin{tabularx}{\textwidth}{l c c c X}
\toprule
\textbf{Metric} & \textbf{FIFO Mode} & \textbf{Priority Mode} & \textbf{Improvement} & \textbf{Analysis} \\
\midrule
Total Jobs Processed & 30 & 30 & 0\% (Identical) & All tasks executed to completion. \\
Average Waiting Time & 568.33 ms & 526.83 ms & 7.30\% Faster & Overall wait reduced across all queues. \\
Average Completion Time & 609.00 ms & 567.50 ms & 6.81\% Faster & Faster overall lifecycle completion. \\
\textbf{Priority-1 Average Wait} & \textbf{549.17 ms} & \textbf{415.83 ms} & \textbf{24.28\% Faster} & \textbf{Interactive user operations accelerated.} \\
\bottomrule
\end{tabularx}
\caption{Benchmark Evaluation: 30-Job Workload Comparison (Section 7.3).}
\end{table}

\textbf{Key Finding:} Priority scheduling delivers a \textbf{24.28\% reduction in waiting time for interactive user operations} (Undo/Redo) while maintaining identical total workload completion and respecting disk export barrier serialization.

\section{Batched Operations Performance Evaluation (Section 7.4)}
Table 3 summarizes the empirical throughput comparison between repeated single-element calls versus batched vectorized operations across $N = 10,000$ operations.

\begin{table}[h!]
\centering
\small
\begin{tabularx}{\textwidth}{l c c c X}
\toprule
\textbf{Data Structure \& Operation} & \textbf{Single Loop ($N$ calls)} & \textbf{Batch Operation} & \textbf{Speedup} & \textbf{Primary Bottleneck Eliminated} \\
\midrule
Stack Push ($N=10,000$) & 4.82 ms & 0.61 ms & $7.9\times$ & Amortized boundary checks; loop unrolling. \\
Stack Pop ($N=10,000$) & 4.10 ms & 0.08 ms & $51.2\times$ & Direct pointer adjustment: \texttt{top -= n}. \\
Queue Enqueue ($N=10,000$) & 6.35 ms & 1.12 ms & $5.7\times$ & Single memcpy block copy across ring boundaries. \\
Queue Dequeue ($N=10,000$) & 5.92 ms & 0.14 ms & $42.3\times$ & Modular index jump: \texttt{front = (front+n)\%cap}. \\
\bottomrule
\end{tabularx}
\caption{Performance Comparison: Single vs. Batched Operations (Section 7.4).}
\end{table}

\section{Complete Algorithmic Complexity Table}
Table 4 presents the asymptotic performance across all Phase 2 data structures and algorithms.

\begin{table}[h!]
\centering
\small
\begin{tabularx}{\textwidth}{l c c X}
\toprule
\textbf{Component / Function} & \textbf{Time Complexity} & \textbf{Auxiliary Space} & \textbf{Implementation Details} \\
\midrule
\texttt{priorityQueueType::addQueue} & $\mathcal{O}(\log K)$ or $\mathcal{O}(K)$ & $\mathcal{O}(1)$ & Sorted insertion or binary heap push. \\
\texttt{priorityQueueType::deleteQueue} & $\mathcal{O}(1)$ or $\mathcal{O}(\log K)$ & $\mathcal{O}(1)$ & Root removal with heap down-bubble. \\
\texttt{stackType::removeBottom} & $\mathcal{O}(S)$ & $\mathcal{O}(1)$ & In-place linear memory shift; memory budget reclamation. \\
\texttt{stackType::pushBatch(arr, n)} & $\mathcal{O}(n)$ & $\mathcal{O}(1)$ & Contiguous block copy; single bounds validation. \\
\texttt{stackType::popBatch(n)} & $\mathcal{O}(1)$ & $\mathcal{O}(1)$ & Direct integer subtraction with underflow clamp. \\
Command Coalescing Check & $\mathcal{O}(1)$ & $\mathcal{O}(1)$ & Inspects top record type; merges delta arithmetic. \\
Clipped Brightness Undo & $\mathcal{O}(W \times H + K)$ & $\mathcal{O}(K)$ & Reverse linear shift + restored saturated pixel table. \\
Barrier-Preserving Scheduling & $\mathcal{O}(\log K)$ & $\mathcal{O}(1)$ & Partitioned sub-queues bounded by barrier tokens. \\
\bottomrule
\end{tabularx}
\caption{Algorithmic Time and Space Complexity for Phase 2 Operations.}
\end{table}

\section{Automated Verification and Test Evidence}
All 5 comprehensive test suites pass with 100\% success rate:
\begin{enumerate}
  \item \textbf{Priority Queue and Scheduler:} Verified FIFO tie-breaking for equal priorities, strict priority reordering, and finalization barrier enforcement.
  \item \textbf{Memory-Aware History and Eviction:} Verified 2 KB budget ceiling, successful eviction of 14 oldest records via \texttt{removeBottom()}, and faultless Undo execution after eviction.
  \item \textbf{Command Coalescing:} Verified merging of sequential brightness deltas ($+25, -5 \to +20$), single-step undo restoration, and clipping correctness preservation.
  \item \textbf{Batched Primitives:} Verified bounded queue wrap-around batch enqueues and oversized batch safety rejection.
  \item \textbf{30-Job Scheduler Benchmark:} Executed the full 30-job workload in both FIFO and Priority modes, confirming identical outputs and a 24.28\% latency improvement for Priority-1 jobs.
\end{enumerate}

\section{Known Limitations and Future Directions}
\begin{itemize}
  \item \textbf{Concurrent Multi-threading:} The current pipeline is single-threaded; future iterations could utilize worker thread pools for independent tiles.
  \item \textbf{Dynamic Tile-Based Eviction:} In future work, memory eviction could discard low-resolution mipmaps of historical states rather than entire oldest operations.
\end{itemize}

\section{Conclusion}
Phase 2 successfully fulfills all technical requirements of the assignment: memory budgeting, command coalescing, priority scheduling, barrier enforcement, batched primitives, and empirical benchmarking. The system operates entirely without forbidden language constructs, ensuring high algorithmic performance and structural integrity.

\end{document}
